% 한국어판: paper/tmlr_ko/build_korean.py 가 생성한다. 직접 고치지 않는다.
% Generated by paper/tmlr/build_assets.py; do not edit by hand.
% input reports/painter_tmlr_diagnostics_v1/analysis.json sha256 76792fa35f90e3d8fec1a75b0eec27d69cac14ed875fce8b64096b8e2525a4d4
% input reports/painter_specificity_v2/psv2-20260911/models.csv sha256 eed121f77626d8a37235b5bb19e370e3fbd7b6bae44b559d1b866e27e9cac41c
\begin{table}[t]
\caption{31개 특징에서의 추가 진단. 반복 잡음: 두 반복의 중심화한 대조값 차이 제곱의 절반을 $H$에 대한 상대값으로 나타낸 것. 보정된 $H$: $H$의 유한 표본 편향(6.7\%)을 제거한 뒤의 $\beta$와 $D$. $\beta$, $Q$, $D-1$이 모두 같은 배율로 조정되므로 순서는 바뀌지 않는다. 중심점 근접도 이득: 일반 지시문에서 화가 지정 지시문으로 바뀔 때 생성 평균에서 각 화가의 참조 평균까지의 제곱 거리가 평균적으로 줄어드는 양으로, 공통 부분과 이름 간 부분으로 정확히 분해된다(부록~\ref{app:estimators}).}
\label{tab:robustness}
\vspace{2pt}
\centering\small
\begin{tabular}{@{}lrrrrrr@{}}
\toprule
 & 반복 잡음 & \multicolumn{2}{c}{보정된 $H$} & \multicolumn{3}{c}{중심점 근접도 이득}\\
\cmidrule(lr){3-4}\cmidrule(l){5-7}
구성 & $/H$ & $\beta$ & $D$ & 전체 & 공통 & 이름 간\\
\midrule
GPT Image 1 & 2.02 & 1.006 & 1.614 & 17.490 & 17.844 & $-$0.354\\
GPT Image 2 & 0.96 & 1.071 & 1.242 & 5.094 & 5.159 & $-$0.065\\
Flare & 0.42 & 0.827 & 1.791 & 0.871 & 1.585 & $-$0.715\\
Sunburst & 0.58 & 0.883 & 1.688 & 4.500 & 4.997 & $-$0.498\\
Nano Banana 2 & 2.60 & 0.474 & 1.117 & 7.423 & 7.013 & 0.410\\
FLUX.2 Max & 1.67 & 0.504 & 0.787 & 10.529 & 10.176 & 0.353\\
\bottomrule
\end{tabular}
\end{table}
